\section{Static-Message Unforgeability of the Derandomized Scheme} \label{app-der-static}

This appendix states and proves the polynomial-assumption result announced in Section~\ref{sec-derand}. The scheme is that of Construction~\ref{constr-der}, with or without the pre-hash; we write it without, since the pre-hash plays no role when no guessing is needed.


\begin{theorem} \label{thm-der-static}
Let $\GG$ be a group of prime order $p$ with generator $g$, and let $f\colon\GG\to\Zp$ be efficient with $\mathsf{Size}_f(0)=0$. Let $\advA=(\advA_0,\advA_1)$ be an adversary that runs in time $t_{\advA}$ and declares $q_s$ signing messages, where $q_s$ is not bounded by any parameter of the scheme. There are an adversary $\advB_{\mathsf{main}}$ against $\CDL[\GG,g,f]$, an adversary $\advB_{\mathsf{trig}}$ against DL, a distinguisher $\advD_0$ against $\PRF_0$, and, for $j\in[q_s]$ and $\beta\in\{0,1\}$, distinguishers $\advD^{\mathsf{io}}_{j,\beta}$ against $\iO$ and $\advD^{\mathsf{prf}}_{j,\beta,i}$ against $\PRF_i$ for $i\in[4]$, together with $\advD^{\mathsf{io}}_{\mathsf{trig}}$ and $\advD^{\mathsf{prf}}_{\mathsf{trig}}$, each running in time $t_{\advA}+\poly(\secp,q_s)$, such that, writing
$$
\epsilon_j
:=
\sum_{\beta\in\{0,1\}}
\Big(
\AdvIO{}{\advD^{\mathsf{io}}_{j,\beta}}
+\sum_{i=1}^{4}\AdvPRF{\PRF_i}{\advD^{\mathsf{prf}}_{j,\beta,i}}
\Big),
$$
we have
$$
\begin{aligned}
\AdvEUFSTATIC{\SchDScheme[\GG,\HF^{\mathsf{der}}]}{\advA}
\le\ &\AdvPRF{\PRF_0}{\advD_0}
+\AdvCDL{\GG,g,f}{\advB_{\mathsf{main}}} \\
&+\AdvDL{\GG,g}{\advB_{\mathsf{trig}}}
+\AdvIO{}{\advD^{\mathsf{io}}_{\mathsf{trig}}}
+\AdvPRF{\PRF_4}{\advD^{\mathsf{prf}}_{\mathsf{trig}}} \\
&+\sum_{j=1}^{q_s}\epsilon_j
+\frac{2q_s}{p} \;.
\end{aligned}
$$
\end{theorem}

\begin{corollary} \label{cor-der-static}
If $\CDL[\GG,g,f]$ holds, $\PRF_0,\ldots,\PRF_4$ and $\iO$ are secure, and $p$ is super-polynomial, then $\SchDScheme[\GG,\HF^{\mathsf{der}}]$ is EUF-static secure against every PPT adversary.
\end{corollary}

\paragraph*{Discussion.} Corollary~\ref{cor-der-static} is ordinary static security of one fixed scheme against every PPT adversary, with no query parameter anywhere and with every assumption polynomially hard. Compared with Corollary~\ref{cor-uf} it removes the bound $Q$, and compared with Theorem~\ref{thm-adaptive} it keeps the signature format of Schnorr. What it does not do is make the queries adaptive; Section~\ref{sec-derand-adaptive} treats that.

Two features of the static game do the work. The declared messages are fixed before key generation, so the reduction knows every slot it will need to program and guesses nothing; this is what keeps the loss polynomial. And the trigger branch recomputes programmed points instead of storing them, so no table appears in the published circuit; this is what removes $Q$.

\begin{proof}[of Theorem~\ref{thm-der-static}]
Let $\advA_0$ declare $m_1,\ldots,m_{q_s}$ and $m^*\notin\{m_1,\ldots,m_{q_s}\}$. Signing is deterministic, so a repeated declared message receives the same signature; we answer repeats by replaying and index the hybrid by the distinct declared messages, which we continue to write $m_1,\ldots,m_{q_s}$.

\heading{Game $\Gm_0$ (uniform nonces).} The game replaces each $r_i=\PRFEv_0(k_0,m_i)$ by an independent uniform $r_i\getsr\Zp$. The key $k_0$ is used nowhere else, and the messages are declared before key generation, so a single distinguisher $\advD_0$ bounds the change by $\AdvPRF{\PRF_0}{\advD_0}$. From here on the honest nonces are uniform, as in the randomized scheme.

\heading{Games $\Gm_1,\ldots,\Gm_{q_s}$ (trigger signing).} For $0\le i\le q_s$, let $\Gm_i$ answer the first $i$ declared messages by the trigger procedure, returning $(W_{m_j},s_{m_j})$, and the remaining ones honestly. Every game publishes the honest circuit $\CircDer$ of Figure~\ref{fig-circuit-derand}; hardwired circuits occur only inside the reductions between adjacent games.

Fix $j$ and compare $\Gm_{j-1}$ with $\Gm_j$. Both endpoints first introduce the event $\mathsf{Bad}_{\mathsf{eq}} := [R_j = W_{m_j}]$ and output $0$ when it occurs, where $R_j=g^{r_j}$ is the honest nonce point. Conditioned on the view before the change, $R_j$ is uniform and independent of $W_{m_j}$, so the two aborts cost at most $2/p$.

The rest of the step is the chain of Appendix~\ref{app-adaptive}, with the slot $u=(d,\salt)$ there replaced by the message $m_j$ here, and with one simplification: the reduction is given $m_j$ by $\advA_0$ before key generation, so it does not guess. Concretely, it sets
$$
t=m_j,\qquad
R=g^{r_j},\qquad
t_R=\enc{R,m_j},\qquad
W=g^{s_{m_j}}X^{-c_{m_j}},\qquad
t_W=\enc{W,m_j},
$$
publishes an obfuscation of the analogue of the two-point circuit $C_j$ of Figure~\ref{fig-circuit-query}, with the slot input $(d,\salt)$ there replaced by the message $m$ and the salt input dropped, punctured at $T=\{t_R,t_W\}$ and at $t$, and runs the eight steps $\mathsf{H}_0$ through $\mathsf{H}_8$: replace $\CircDer$ by $C_j$ at the true values, make the $\PRF_1$ and $\PRF_2$ values on $T$ uniform, replace the main-branch hash value by a uniform value and discard the masks, change variables so the honest record is $(g^{s_0}X^{-v},v,s_0)$ for uniform $v,s_0$, make the $\PRF_3$ and $\PRF_4$ values at $t$ uniform, exchange the two records by Lemma~\ref{lem-swap}, and reverse. Each of the two directions invokes $\iO$ once and the four puncturable PRFs once, giving $\epsilon_j$, and the puncture sets are fixed by data available before key generation, so Definitions~\ref{def-pprf} and~\ref{def-io} apply with polynomial-time samplers. Hence
$$
\big|\Pr[\Gm_{j-1}\Rightarrow1]-\Pr[\Gm_j\Rightarrow1]\big|
\le \epsilon_j+\frac{2}{p} \;.
$$

In $\Gm_{q_s}$ every signature is a trigger signature, and computing $W_{m}$ and $s_{m}$ uses $X$ and the PRF keys but not $x$.

\heading{A main-branch forgery.} Let $E_{\mathsf{main}}$ be the event that $\advA_1$ returns a valid forgery on $m^*$ whose hash value comes from the main branch. Adversary $\advB_{\mathsf{main}}$ receives a CDL challenge $h$, sets $X\gets h$, samples the four PRF keys, publishes the obfuscated circuit, and answers the declared messages by the trigger procedure. Given a valid forgery $(R^*,s^*)$ it computes $a^*,b^*$ from $k_1,k_2$ at $\enc{R^*,m^*}$ and returns
$$
\big(R^*X^{a^*}g^{b^*},\ (s^*+b^*)\bmod p\big),
$$
which wins the CDL game by the computation of Section~\ref{sec-io-constr}, using that $f$ has no zeros. Hence $\Pr[E_{\mathsf{main}}]\le\AdvCDL{\GG,g,f}{\advB_{\mathsf{main}}}$.

\heading{A trigger-branch forgery.} Let $E_{\mathsf{trig}}$ be the event that the forgery is valid and its hash value comes from the trigger branch, so that $R^*=W_{m^*}$. Here the static setting removes the guessing of Appendix~\ref{app-adaptive}: the forgery message $m^*$ is declared by $\advA_0$ before key generation, so the reduction knows the slot it must program.

Adversary $\advB_{\mathsf{trig}}$ receives a DL challenge $Y$, samples $x\getsr\Zp$ and sets $X\gets g^x$, punctures $k_4$ at $m^*$, and publishes an obfuscation of the analogue of the circuit $C_{\mathsf{forg}}$ of Figure~\ref{fig-circuit-trigger-forgery}, with the slot input retyped to the message in the same way, with $\widehat t=m^*$ and $Z=Y$. With $Z=g^{\widehat s}$ for $\widehat s=\PRFEv_4(k_4,m^*)$ this circuit computes the same function as $\CircDer$, so $\iO$ justifies the switch, and puncturable-PRF security at $m^*$ replaces $\widehat s$ by a uniform value, after which $Z=Y$ has the right distribution. Signing uses the punctured key and never needs slot $m^*$, since $m^*$ is not among the declared messages. If the forgery takes the trigger branch at $m^*$ then
$$
g^{s^*}=\big(YX^{-c_{m^*}}\big)X^{c_{m^*}}=Y,
$$
so $\advB_{\mathsf{trig}}$ returns $s^*$ as the discrete logarithm of $Y$. Hence
$$
\Pr[E_{\mathsf{trig}}]\le
\AdvDL{\GG,g}{\advB_{\mathsf{trig}}}
+\AdvIO{}{\advD^{\mathsf{io}}_{\mathsf{trig}}}
+\AdvPRF{\PRF_4}{\advD^{\mathsf{prf}}_{\mathsf{trig}}} \;.
$$

\heading{Conclusion.} A valid forgery takes one of the two branches, so collecting the terms gives the stated bound. Every reduction runs in time $t_{\advA}+\poly(\secp,q_s)$, and no factor depending on the message space or on a guess appears.
\end{proof}

